\documentclass{article}
\usepackage{amsmath,amssymb,amsthm}
\newtheorem{theorem}{Theorem}
\begin{document}

\begin{theorem}
Let \(n=8\). On an interval where the eight operators
\[
srx,\ sxp,\ cxp,\ crx,\ urx,\ sawup,\ uxp,\ sawdown
\]
are defined and nonzero (for example \(0<x<\pi/2\)), set
\[
a_1=srx,\quad a_2=sxp,\quad a_3=cxp,\quad a_4=crx,\quad
a_5=urx,\quad a_6=sawup,\quad a_7=uxp,\quad a_8=sawdown.
\]
Define the \(8\times 8\) matrices
\[
X_i=a_i E_{i,i+1}\qquad (i=1,\dots,8),
\]
where the indices are taken modulo \(8\), so \(E_{8,1}\) denotes the \((8,1)\)-entry. Then the Lie algebra \(\mathfrak{g}=\langle X_1,\dots,X_8\rangle\) generated by these matrices is \(\mathfrak{sl}(8)\). Consequently,
\[
\dim\mathfrak{g}=63,\qquad \operatorname{rank}\mathfrak{g}=7.
\]
\end{theorem}

\begin{proof}
Each \(a_i\) is a unit in the coefficient field. For any \(i\),
\[
[X_i,X_{i+1}]
=[a_iE_{i,i+1},\,a_{i+1}E_{i+1,i+2}]
=a_i a_{i+1}E_{i,i+2}.
\]
Since \(a_i a_{i+1}\) is a unit, we may divide by it and obtain
\[
E_{i,i+2}\in\mathfrak{g}.
\]

Now suppose inductively that \(E_{i,i+k}\in\mathfrak{g}\) for some \(1\le k\le 6\). Then
\[
[E_{i,i+k},X_{i+k}]
=[E_{i,i+k},\,a_{i+k}E_{i+k,i+k+1}]
=a_{i+k}E_{i,i+k+1}.
\]
Dividing by the unit \(a_{i+k}\), we get
\[
E_{i,i+k+1}\in\mathfrak{g}.
\]
Thus, by induction, every elementary matrix \(E_{i,j}\) with \(i\ne j\) lies in \(\mathfrak{g}\). All indices are modulo \(8\); a path of length \(k\) around the directed cycle corresponds to the product of the intervening \(a\)'s, which is a unit.

In particular, for every \(i\ne j\), both \(E_{i,j}\) and \(E_{j,i}\) belong to \(\mathfrak{g}\). Hence
\[
[E_{i,j},E_{j,i}]=E_{i,i}-E_{j,j}\in\mathfrak{g}.
\]
The matrices \(E_{i,i}-E_{j,j}\) span the diagonal traceless subspace
\[
\mathfrak{h}=\left\{\operatorname{diag}(d_1,\dots,d_8):\sum_{i=1}^8 d_i=0\right\}.
\]
The off-diagonal elementary matrices \(E_{i,j}\) \((i\ne j)\) together with \(\mathfrak{h}\) span the full trace-zero subspace
\[
\mathfrak{sl}(8)=\{M\in\mathfrak{gl}(8):\operatorname{tr}M=0\}.
\]
Since all these elements have been shown to lie in \(\mathfrak{g}\), we conclude
\[
\mathfrak{g}=\mathfrak{sl}(8).
\]

Finally,
\[
\dim\mathfrak{sl}(8)=8^2-1=63,\qquad
\operatorname{rank}\mathfrak{sl}(8)=8-1=7,
\]
and the root system is of type \(A_7\).
\end{proof}

\end{document}
